% Generated draft: requires mathematical and source review.
\documentclass[11pt]{article}
\usepackage[margin=25mm]{geometry}
\usepackage{fontspec}
\setmainfont{Latin Modern Roman}
\usepackage{amsmath,amssymb,mathtools,unicode-math}
\setmathfont{Latin Modern Math}
\usepackage{enumitem,needspace,xurl,hyperref}
\setlength{\parindent}{0pt}
\setlength{\parskip}{5pt}
\setlength{\emergencystretch}{3em}
\newcommand{\N}{\mathbb{N}}
\newcommand{\Z}{\mathbb{Z}}
\newcommand{\Q}{\mathbb{Q}}
\newcommand{\R}{\mathbb{R}}
\newcommand{\C}{\mathbb{C}}
\newcommand{\F}{\mathbb{F}}
\newcommand{\A}{\mathbb{A}}
\newcommand{\PP}{\mathbb{P}}
\newcommand{\E}{\mathbb{E}}
\newcommand{\eps}{\varepsilon}
\newcommand{\dd}{\,\mathrm{d}}
\begin{document}
\section*{UnsolvedMath ID: 20000814}
\textbf{Source problem code:} catalogue entry
\subsection{Definitions and mathematical statement}
Let $k$ be algebraically closed of characteristic zero. In the flat-specialization interpretation, let $\mathcal C\subset\PP_R^3$ be a flat projective family of curves over a DVR with residue field $k$, with smooth geometric generic fiber a complete intersection of surfaces of degrees $1\le a\le b$ and smooth special fiber $C$. The source asks whether $C$ must also be a complete intersection. Flatness and a trait base make precise the original wording ``family,'' ``general'' and ``special''; they are not explicit in the original sentence.

The attached aid states that the special fiber retains
\[
\deg C=ab,\qquad p_a(C)=1+\frac{ab(a+b-4)}2,\qquad
\omega_C\simeq\mathcal O_C(a+b-4),
\]
and satisfies $H^0(\PP_k^3,\mathcal I_C(a))\ne0$, so it lies on a degree-$a$ surface. Define its Hartshorne--Rao module by
\[
M_C=\bigoplus_{m\in\Z}H^1(\PP_k^3,\mathcal I_C(m)).
\]
Within this smooth subcanonical setting, the aid reports that $C$ is a complete intersection if and only if $M_C=0$. It also reports that a smooth containing surface of degree $a$ forces $C$ to be a complete intersection.

A direct analysis of planes and quadrics of every rank is said to prove the original assertion for $a\le2$. Consequently any counterexample to this flat characteristic-zero question must have $a\ge3$, nonzero Rao module, the displayed subcanonical equality, and every degree-$a$ containing surface singular. These are simultaneous necessary conditions from the aid, not a construction or a proof of the assertion for all $a$.

The original remark records negative answers if the characteristic is positive, if the ambient projective dimension is greater than $3$, or if the special fiber is allowed to be nonsmooth, and supplies Ellia--Hartshorne as the reference for those qualifications.
\subsection{Short English statement}
Must a smooth characteristic-zero limit of smooth complete-intersection space curves remain a complete intersection? The supplied aid proves the low first-degree cases and restricts any counterexample to nonzero Rao module and singular minimal-degree containing surfaces, while preserving the original counterexamples outside the stated hypotheses.
\subsection{Sources}
\begin{itemize}
\item Local UnsolvedMath registry, permanent ID $20000814$, source code catalogue entry; canonical \texttt{aim-arithmetic-geometry-notes.json}, \texttt{notes[59]}.
\item American Institute of Mathematics, workshop ``Components of Hilbert schemes,'' problem 13: \url{https://aimath.org/WWN/hilbertschemes/hilbertschemes.pdf}. No PDF inspection is claimed here.
\item Original remark: P. Ellia and R. Hartshorne, \emph{Smooth specializations of space curves: questions and examples}, in \emph{Commutative algebra and algebraic geometry} (Ferrara), pp. 53--79, Lecture Notes in Pure and Applied Mathematics 206, Dekker, New York, 1999.
\item Attached research attempt 1: partial result, recorded run status open, rigor proved, novelty confidence low; registry status partially solved, classification PARTIAL-PROGRESS. The reductions and low-degree claim are machine generated and not peer reviewed; independently verify before citation.
\end{itemize}
\end{document}
